\section{Rigorous computation and reproducibility}\label{sec:computation}

This section describes what the computer-assisted parts of the proof consist of, what has to be trusted, and how
a reader can check them.

\subsection{Arithmetic and trust base}\label{sec:computation-arith}

Three kinds of computation enter the certificates.
\begin{enumerate}[label=(\alph*),leftmargin=*]
\item \emph{Exact rational arithmetic} (Python's \texttt{fractions}): the final inequalities of every
  certificate, namely the Perron conditions $Nv\le\theta v$ with $\theta<1$ of
  \cref{lem:perron}, the Newton--Kantorovich conditions of \cref{sec:roots-nk}, the covers of the contours and of
  the region of \cref{prop:Kinj}, the deflation polynomial $q$, and the fact F1.
\item \emph{Ball arithmetic} with Arb~\cite{Arb}, through \texttt{python-flint}: the symbol bounds of
  \cref{lem:F3}, and the slack between grid points.
\item \emph{A posteriori bounds on floating-point computations.} Large linear algebra (dimension $14016$ for the
  finite reference, windows of dimension up to $11200$) is done in IEEE double precision. Floating-point results
  are used only as approximate inverses, solvers or test vectors, and every quantity that enters a certificate is
  bounded rigorously: residuals are computed with entrywise rounding bounds of the standard form
  $\gamma_n\||M||X|\|$~\cite{Higham2002}; norms of inverses are bounded by verified Loewner inequalities
  $t^2MM^*-I\succeq0$, whose positive definiteness is verified by Rump's method~\cite{Rump2006} applied to the
  realification; and the coupling factors are certified as in \cref{sec:counting-reference}.
\end{enumerate}
The trust base is therefore: the correctness of the bounds published in~\cite{RT2019}, which we use as stated;
IEEE~754 arithmetic with rounding to nearest, as assumed by the a posteriori bounds; Arb; and the Python code of the
verification program, which is short for each check and is available for inspection. No part of the final chain
is formalized in a proof assistant.

Because the floating-point entries depend on the BLAS library and on the number of threads, re-running a
certificate on another machine reproduces the rational $\theta$ only up to $\sim10^{-16}$, not exactly; the
criterion for a successful reproduction is $\theta<1$ with a difference at the scale of rounding.

\subsection{The certificates}\label{sec:computation-certs}

\Cref{tab:certs} lists the computed certificates, the results they establish, and their approximate cost. The
computations were run on five commodity computers, with four to six cores and $7$--$31$\,GB of memory; the whole
chain takes a few days on three or four such machines. The largest single computations are the Schur
decomposition of the finite reference and the Loewner bounds on its resolvent (hours per contour centre), and the
tails with $26$ dense windows (three to six hours per centre).

\begin{table}[ht]
\centering
\small
\begin{tabular}{@{}p{0.29\textwidth}p{0.35\textwidth}p{0.27\textwidth}@{}}
\toprule
certificate & establishes & cost \\
\midrule
symbols of $B$ and $B^\sharp$ (Arb, grids $2048\times4096$ and $1024\times2048$) & \cref{lem:F3}; \cref{thm:counts}~(i) & hours, 4 processes \\
radius recovery and transfer & \cref{lem:Lreal,lem:recovery} & minutes \\
finite reference: matrix, Schur, residual & \cref{lem:schur} & hours \\
level $Z$ on $174$ discs (Loewner, Schur solves) & \cref{thm:counts}~(ii),(iii) & $\sim$9 min per point \\
low-rank coupling factors ($13$ centres) & \cref{sec:counting-reference} & $\sim$1 h per centre \\
tails, $26$ windows ($13$ centres) & \cref{sec:counting-reference} & 3--6 h per centre \\
windows inside the bands & \cref{lem:windows} & hours \\
bordered NK at $\mu,\sK,\sphys,\sB$ & \cref{prop:nk,prop:mu} & 3--6 h each \\
constraint witnesses & \cref{prop:witness} & minutes \\
$154$ tiles and Rouch\'e for $K$ & \cref{prop:Kinj} & $\sim$1 day, 3 machines \\
F1 & \cref{lem:F1} & seconds \\
\bottomrule
\end{tabular}
\caption{Computed certificates. File names and sha256 hashes are listed in \cref{app:certificates}.}
\label{tab:certs}
\end{table}

\subsection{The verification program}\label{sec:computation-verify}

The repository contains a verification program with two levels.

\emph{Level 1} re-checks, in exact rational arithmetic and with the code of the repository, the final inequalities
of the certificates from the stored certificate files (small JSON files):
\begin{itemize}[leftmargin=*]
\item the three-level Perron bound at each of the $175$ points of the $22$ disc files of the two bands ($174$ discs of
  positive radius, and one entry of radius zero that is not used in the cover);
\item the covers of $\partial\Omega_A$ and $\partial\Omega_B$;
\item the finite counts;
\item the window bounds;
\item the four Newton--Kantorovich certificates, the witnesses and the identification of the gauge root;
\item the $154$ tiles and their cover of the region of \cref{prop:Kinj} outside the disc $|s-0.401|<0.05$ of the
  Rouch\'e argument for $K$;
\item the deflation polynomial $q$ of \cref{sec:counting-rouche};
\item the radius recovery and transfer bounds, the bound $\beta_+$ on the background part, and F1.
\end{itemize}
Some certificates enter level 1 only through their stored values, and are recomputed only at level 2: the symbol
bounds and the free parts (level 1 recombines the stored bands), and the Perron bounds of the $16$ tiles of the
Rouch\'e argument for $K$.
It takes about two hours on four cores. All checks pass: the worst Perron root on the discs is $0.999074$, with
differences to the stored values $\le7.8\cdot10^{-16}$, and the Newton--Kantorovich data, witnesses and tiles are
reproduced identically.

\emph{Level 2} recomputes the intermediate certificates from the data of~\cite{RT2019}. As a sample, we recomputed
from scratch on a different machine:
\begin{itemize}[leftmargin=*]
\item the reference matrix, identical bit for bit to the stored one;
\item the Schur residual bound, identical;
\item the tail of the contour centre $3.0$ ($26$ windows and the far level, $5.3$ hours), identical in every field.
\end{itemize}

\subsection{Data and code availability}\label{sec:computation-data}

The code, the certificate files, the verification program and a reproduction guide are available at
\url{https://github.com/edivaldo-amaral/choptuik-espectro} (MIT licence)~\cite{Repo2026}. The large data files, namely the reference
matrix, its Schur factors and the coupling factors ($16$\,GiB), are deposited at Zenodo, DOI
\href{https://doi.org/10.5281/zenodo.23222363}{10.5281/zenodo.23222363} (CC BY 4.0)~\cite{Zenodo2026}, with sha256 hashes that are
also recorded in the certificate files. The data of~\cite{RT2019} are not redistributed; a script downloads them
from the arXiv source of~\cite{RT2019} and checks their sha256 hashes.
\todo{Repository and Zenodo record become public with the preprint.}

\subsection{Independent AI-assisted checks}\label{sec:computation-reviews}

Before being used, the certificates and the lemmas were checked by AI agents independent of those that produced the
work, with the exceptions listed below. These checks are not reviews by human experts. The protocol was designed so
that a check produces verifiable evidence rather than an opinion:
\begin{itemize}[leftmargin=*]
\item the checking agent ran on a different model from the one that produced the work, with a clean context;
\item it was given the precisely stated claims, the derivations under review, and the code and data, but not
  the author's conclusions or earlier audits;
\item it had to deliver one artifact per claim: an independent derivation, an independent implementation, or a
  direct computation of the true quantity to compare with the bound;
\item a snapshot with sha256 manifest was taken before each check and compared after it.
\end{itemize}
There were ten such checks:
\begin{itemize}[leftmargin=*]
\item the tiles of \cref{prop:Kinj} (an earlier check of them with another tool was superficial and is not
  counted);
\item the certificates at $\sK$ (with the Rouch\'e argument for $K$) and at $\mu$;
\item the Rouch\'e arguments of the two bands, including the certificate at $\sB$;
\item the structural lemmas and the chain of the proof;
\item the gauge lemma, the forward reconstruction and the converse;
\item the modelling assumptions.
\end{itemize}
Two parts have not been checked on their own. The Newton--Kantorovich certificate at $\sphys$ uses the same code as
the one at $\sB$, which was checked. The symbol computations of \cref{lem:F3} were examined only in their use, not
recomputed. None of the checks overturned a result. All of them found gaps, mostly of traceability or of statement,
and some of derivation: for instance, one found that an early version of the tiles used the dyadic approximation of
$\mu$ instead of the true value, and another that the deflation error outside the finite box had been omitted. All
were corrected before the results were used; the reports and artifacts are in the repository.

\subsection{Use of AI assistants}\label{sec:computation-ai}

This work was done with extensive use of AI assistants, between September and October 2026. We describe their role
as required by the journal's policy.

\emph{Scope.} The assistants wrote most of the code, carried out most of the derivations and computations, prepared
the documentation, and drafted this manuscript. Their contribution is substantial and not limited to language
editing.

\emph{The author's role.} The author posed the problem, the unstable spectrum of the solution of~\cite{RT2019}, which
its authors left open; directed the project and chose among the approaches; ran the computations on the author's own
and on borrowed machines; checked the arguments and the results; and takes full responsibility for the content of
this paper.

\emph{Tools.} The main assistant was Claude (Anthropic), used through the Claude Code environment, with the models
\texttt{claude-opus-5} and later \texttt{claude-opus-5-5}. OpenAI's Codex agent was used in the first exploratory
session and during a three-day interval, with the models \texttt{gpt-5.6-luna}, \texttt{gpt-5.6-sol},
\texttt{gpt-6-astra} and \texttt{gpt-reserve}; everything produced in that interval was audited and re-derived or
checked before use. The checks of \cref{sec:computation-reviews} were done by Claude agents running
\texttt{claude-fable-5-1}, a model different from the one that produced the work. At the author's request, a draft of the
manuscript was also read by ChatGPT (OpenAI, model Sol 6.1); its comments were checked against the certificates and
the sources before being incorporated. The model identifiers are those recorded in the session logs; a dated account of the sessions is in the file \texttt{paper/AI\_USE.md} of the
repository.

\emph{Limits of the verification.} The analytic arguments are given in this paper; some routine estimates appear in
sketched form, with their constants checked by code in the repository (\cref{app:lemmas}). Every computed bound is
produced by code that can be inspected and re-run from the certificate files and data, and all of them have been
re-checked by the verification program (level~1 of \cref{sec:computation-verify}); a sample of the intermediate
certificates has been recomputed from the data of~\cite{RT2019} (level~2). The lemmas and most certificates were also
examined by the checks of \cref{sec:computation-reviews}, with the exceptions listed there. None of this replaces
peer review; the aim is that trust in the result should not depend on the process that produced it.
